% Legislative Findings and Evidentiary Record. Built by the future pass from the
% instructions below. Do not process now. This is Theme B of the prompt: the
% recorded evidence the bill rests on. Cast each finding in the numbered-finding
% legislative style ("The Congress finds the following:"), then support it from
% the named records. Keep the dense numbers in the evidence table.

[Open with the standard findings preamble, then enumerate the findings in three
grouped subsections. Each finding is one or two sentences of legislative
declaration supported by a sentence that points (by table reference, not inline
path) to the recorded evidence. Build one evidence table (Finding, Source record,
Headline metric) using the left-aligned ragged-right columns, with the parent
paths given once at the top of each block. Do not narrate the table row by row.]

[Group A, LLMs are appropriate for VVUQ code verification (two author studies).
Finding A1, from VVUQ-01: an autonomous large language model agent built,
executed, and documented a VVUQ pipeline under strict assurance, and the gate, not
the generation, was the decision-bearing step (the pipeline emitted a complete
deliverable in under a millisecond while the gate accepted only one of six
candidates, blocking five and escalating one, with the verification fraction held
at exactly 1.0 and 51 of 51 tests passing). Finding A2, from VVUQ-02: the same
priority held for an embodied surgical humanoid, where the assurance suite carried
64 of 172 tests, more than a third of the budget, evidence in itself that
verification outweighs generation. Source from
\texttt{VVUQ-01/final-paper/sections/} (results, discussion, methods) and
\texttt{VVUQ-01/execution/README.md}, and from
\texttt{VVUQ-02/final-paper/sections/} and \texttt{VVUQ-02/execution/README.md}.
Cite \texttt{kawchak2026vvuq01} and \texttt{kawchak2026vvuq02}.]

[Group B, Practical VVUQ code generation and execution against External Standards
(the substantial external-standard incorporation). This group is required to be
comprehensive about the External Standards, which are the heart of the VVUQ-02
contribution. State the finding that VVUQ code generation and execution are
practical today and, crucially, that every gate was bound to published consensus
standards already used in real life rather than to ad hoc thresholds, which is
what makes the assurance argument traceable and defensible to a regulator. Pull
the standards corpus and the per-gate bindings from
\texttt{VVUQ-02/codegen/docs/standards\_map.md} and the per-standard input
summaries under \texttt{VVUQ-02/inputs/standards/}, and the methodology from
\texttt{vvuq\_methodology.md} and \texttt{vvuq\_gate\_spec.md}. Enumerate the
fourteen standards and two clinical baselines (ASME V\&V 40-2018, NASA-STD-7009A,
IEC 80601-2-77, IEC 60601-1, ISO/TS 15066, ISO 13482, ISO 10218-1, ISO 9283,
IEC 62304, ISO 14971, ISO 13849-1, UL 4600, IEEE 7009, the FDA credibility
guidance, the Dutch cohort, and the Callery Fistula Risk Score) in a table, not
in prose, and cite each: \texttt{asmevv40}, \texttt{nasastd7009a},
\texttt{iec8060127}, \texttt{iec60601}, \texttt{isots15066}, \texttt{iso13482},
\texttt{iso10218}, \texttt{iso9283}, \texttt{iec62304}, \texttt{iso14971},
\texttt{iso13849}, \texttt{ul4600}, \texttt{ieee7009}, \texttt{fda2023credibility}.
Reference, by table, the two featured artifacts processed and executed (the
1790-line four-entrant comparison and the 1001-row non-repetitive humanoid sensor
stream) but reserve their detailed treatment for
\texttt{sections/algorithm\_documentation.tex} and
\texttt{sections/supporting\_documentation.tex}.]

[Group C, Physical AI VVUQ is supported by two author standards and five trial,
sponsor, and mobile simulations. State the finding that the readiness standards
and the simulation record make the robot-patient case concrete. The two author
standards are USL and PSL, sourced from \texttt{psl\_usl\_standards.tex} and the
site documentation package; cite \texttt{kawchak2026national} and
\texttt{kawchak2026sitedocs}. The five simulations, to be listed in the evidence
table with one headline metric each, are: (1) the national-platform 24-hour
multi-patient trial simulation (168 patients, 29 robots, 99.7 percent uptime,
zero harm) from \texttt{site\_establishment.tex}; (2) the single-patient ten-stage
journey (Stage IIIB NSCLC over 1,120 days) from \texttt{patient\_journey.tex};
(3) the fully automated sponsor simulation (288 sponsor decisions across a
24-hour trial), \texttt{kawchak2026sponsor}; (4) the mobile pancreatic-cancer
Unitree H2 surgical humanoid simulation, \texttt{kawchak2026vvuq02}; and (5) the
accelerated patient-prediction study of four extensive simulations,
\texttt{kawchak2026prediction}. Add the federated-learning triple-AI peer review
as supporting infrastructure, \texttt{kawchak2026federated}.]

% Model evidence table (parent paths at top of each block, abbreviated leaves):
% \begin{table}[h]\centering
% \begin{tabularx}{\textwidth}{L{2.0cm} L{5.2cm} Y}
% \toprule
% \textbf{Finding} & \textbf{Source record (abbreviated)} & \textbf{Headline metric} \\
% \midrule
% \multicolumn{3}{l}{\textit{Parent: papers/VVUQ-01/ and papers/VVUQ-02/}} \\
% A1 & VVUQ-01 execution/README.md & 51/51 tests; 1 ACCEPT, 5 BLOCK, 1 ESCALATE \\
% A2 & VVUQ-02 execution/README.md & 172 tests; 64 in the 10-gate suite \\
% B  & VVUQ-02 codegen/docs/standards\_map.md & 14 standards + 2 clinical baselines, per gate \\
% \multicolumn{3}{l}{\textit{Parent: physical-ai .../final\_paper/sections/ and Zenodo}} \\
% C1 & site\_establishment.tex & 168 patients, 29 robots, 99.7\% uptime, 0 harm \\
% C2 & patient\_journey.tex & 1 patient, 10 stages, 1{,}120 days \\
% C3--C5 & sponsor; H2 humanoid; prediction (Zenodo) & 288 decisions; mobile Whipple; 4 simulations \\
% \bottomrule
% \end{tabularx}
% \caption{Evidentiary record underlying the legislative findings.}
% \end{table}
